% !TEX program = xelatex
% 离散数学(2) 期末 cheatsheet —— 2张A4双面(4页) 全定理收录+算法全展开
\documentclass[4.5pt,landscape]{extarticle}
\usepackage[a4paper,margin=0.18cm]{geometry}
\usepackage{ctex}
\usepackage{amsmath,amssymb}
\usepackage{multicol}
\usepackage{enumitem}
\usepackage{xcolor}
\usepackage{titlesec}

\setlength{\columnsep}{3.5pt}
\setlength{\parindent}{0pt}
\setlength{\parskip}{0.5pt}
\setlength{\abovedisplayskip}{0pt}\setlength{\belowdisplayskip}{0pt}
\linespread{0.74}

\definecolor{rd}{RGB}{170,25,25}
\definecolor{bu}{RGB}{25,55,180}
\definecolor{gr}{RGB}{0,105,0}
\newcommand{\rd}[1]{\textcolor{rd}{\textbf{#1}}}
\newcommand{\bu}[1]{\textcolor{bu}{\textbf{#1}}}
\newcommand{\xt}[1]{\textbf{#1}}
% 算法: 用 flushleft + \\ 换行 + \quad 缩进, 不再依赖 tabbing
\newcommand{\alg}[1]{\par\vspace{0pt}{\ttfamily\footnotesize\color{gr}\raggedright #1\par}\vspace{0pt}}
\newcommand{\ST}{\hspace*{0.9em}}
% 打散数学公式拥挤
\newcommand{\gap}{\,}

\titleformat{\section}{\bfseries\color{rd}\fontsize{5pt}{5.5pt}\selectfont}{}{0pt}{}[\vspace{-2pt}\textcolor{rd}{\hrule height 0.25pt}]
\titleformat{\subsection}{\bfseries\color{bu}\fontsize{5pt}{5.5pt}\selectfont}{}{0pt}{}
\titlespacing*{\section}{0pt}{1.2pt}{0.1pt}
\titlespacing*{\subsection}{0pt}{0.5pt}{0pt}
\setlist[itemize]{leftmargin=5pt,topsep=0pt,itemsep=0pt,parsep=0pt,label=\textbullet}

\pagestyle{empty}
\begin{document}
\begin{multicols*}{3}
{\centering\xt{\large 离散数学(Ⅱ) 期末 Cheatsheet}\par\vspace{0pt}}

% ====== 图论基础 ======
\section{图论基础}
$G=(V,E)$, $n=|V|$, $m=|E|$, 默认有限无向简单图。$d(v)$度数(自环贡献2), $\Delta=\max d$, $\delta=\min d$。

\rd{握手定理: 任意图所有顶点度数之和$=2m$; 有向图 $\sum d^+=\sum d^-=m$。}\rd{推论: 任意图奇度顶点的个数必为偶数。}

\rd{非空简单图一定存在度数相同的顶点}。

$K_n$ (无向完全图): $m=\binom n2=n(n-1)/2$; 有向完全图: $m=n(n-1)$, $\Delta^+=\delta^+=n-1$。\xt{$k$-正则图}: 每点度均为$k$, $m=kn/2$。\xt{圈图$C_n$}: 至少3点, 2-正则图, 每个点度均为2。\xt{轮图$W_n$}: 圈$C_{n-1}$内加中心点连所有点, 中心度$n-1$外圈点度3。

\xt{二分图$G=(V_1,V_2,E)$}: $V_1\cap V_2=\emptyset$, 所有边均为$V_1$与$V_2$之间。完全二分图$K_{m,n}$有$mn$条边。二分图可能有多种剖分方式。

\xt{同构$G_1\cong G_2$}: 存在双射$f:V_1\to V_2$保持邻接关系。充要条件: 点数、边数、度序列一致, 且对应回路结构一致(必要条件, 对简单图常足够)。

图运算: 删边$G-e$, 加边$G+e$, 删点$G-v$, 边收缩$G/e$(合并两端点,删$e$), 并$G_1\cup G_2$, 交$G_1\cap G_2$, 差$G_1-G_2$, 笛卡尔积$G_1\times G_2$。

\xt{邻接矩阵$A$}: $a_{ij}=v_i$到$v_j$的边数, 无向对称。\bu{$A^\ell$ 的 $(i,j)$ 元 $=v_i$ 到 $v_j$ 长度为 $\ell$ 的通路数}(归纳法\xt{可达矩阵$P$}: $P=A\vee A^2\vee\cdots\vee A^n$(布尔运算), 或$P_{ij}=1$当$i$可达$j$。Warshall/Floyd: 三重循环$O(n^3)$, 可达$p_{ij}^{(k)}=p_{ij}^{(k-1)}\vee(p_{ik}^{(k-1)}\wedge p_{kj}^{(k-1)})$; 最短路$d_{ij}^{(k)}=\min(d_{ij}^{(k-1)},d_{ik}^{(k-1)}+d_{kj}^{(k-1)})$。

\xt{关联矩阵$B$}(无环有向图): $b_{ij}=1$若$v_i$为$e_j$始点; $=-1$若终点; $=0$不关联。无向图: $b_{ij}=1$关联, $0$不关联(自环贡献2)。性质: 每列和$=0$(有向); 握手定理$\sum_{i,j}b_{ij}=2m$(无向)。邻接表/正向表/逆向表/十字链表是存储优化。

% ====== 道路与回路 ======
\section{道路·回路·连通性}
\xt{通路(Walk)}: 点边交替序列; \xt{迹(Trail)}: 边不重; \xt{初级道路(Path)}: 点不重; \xt{闭通路(Closed Walk)}; \xt{圈(Cycle)}: 闭的初级道路($\geq3$)。长度=所含边数。

\xt{连通}: 无向图任两点间有路。\xt{连通分支数$p(G)$}。\xt{点割集}: 删后连通分支数增加的点集; \xt{割点}: 单点构成点割集。\xt{边割集/桥}: 删后连通分支数增加的边集; 桥是单边边割集。\xt{点连通度$\kappa(G)$}=最小点割基数, 规定$\kappa(K_n)=n-1$, $\kappa($不连通$)=0$。\xt{边连通度$\lambda(G)$}=最小边割基数。\xt{基本关系}: \bu{$\kappa(G)\leq\lambda(G)\leq\delta(G)$}。推论: 简单连通图 $\kappa(G)\leq 2m/n$。

\xt{块(Block)}: 极大的没有割点的连通子图。\rd{定理2.1.3(块判别,$n\geq3$连通图等价)}: (1)$G$是块$\Leftrightarrow$(2)任两点同属某初级回路$\Leftrightarrow$(3)任一点与任一边同属某初级回路$\Leftrightarrow$(4)任两边同属某初级回路$\Leftrightarrow$(5)给定两点$u,v$和一边$e$, 存在含$e$的初级路$P_{uv}$$\Leftrightarrow$(6)任意三个不同顶点存在包含它们的初级路$\Leftrightarrow$(7)任意三个不同顶点存在只含两点不含第三点的路。\xt{意义: 块中任两点都在某个圈上, 块是2-连通的基本单位。}


\subsection{欧拉道路与回路}
\bu{定理2.3.1: 无向连通图$G$存在欧拉回路(每边恰一次)的充要条件是$G$中各顶点的度均为偶数。}
\bu{推论2.3.1: 若无向连通图$G$中只有两个奇度顶点, 则$G$存在欧拉道路(通)且道路端点恰为这两个奇度点。}
定理2.3.2: 连通图有$k$个奇度点, 则$E(G)$可划分成$k/2$条简单道路。推论2.3.2: 图$G$连通且恰有$2k$个奇点, 则$G$可用$k$笔不重复画成且至少$k$笔。
\bu{推论2.3.3: 有向连通图$G$中各节点正度等于负度(入度=出度), 则$G$中存在有向欧拉回路。}
推论2.3.4: 有向连通图恰有两个节点正负度不等(一入比出多1, 一入比出少1), 则存在有向欧拉通路。

\xt{Fleury算法}: 从奇度点(或任点)出发, 每次优先选非桥边走, 走过即删边。朴素$O(m^2)$(每次需判桥)。
\alg{任选起点$v_0$, 栈空; \\ while 存在未访问边: \\\ST 从当前点DFS沿未访问边走, 走过即删边 \\\ST 无法前进时回溯, 将点压入栈 \\ 最终栈中逆序即为欧拉回路(或道路)}

\subsection{哈密顿道路与回路}
哈密顿路/回路: 经过所有顶点恰好一次的初级道路/回路。\xt{实质}: 能将图中所有顶点排在同一圈中。欧拉与H图无包含关系。

\rd{必要条件(Th2.4.1必要)}: 若图$G$有H回路, 则对$V$的任何非空子集$S$, 有 $p(G-S)\leq |S|$, 其中$p(G-S)$是删$S$及关联边后的连通分支数。\xt{(逆否命题: 若存在$S$使$p(G-S)>|S|$, 则$G$无H回路)}
\bu{推论1(道路): 若$G$有H道路, 则 $p(G-S)\leq|S|+1$。}
\bu{推论2: 有割点的图不是哈密顿图}(取$S=\{v\}$割点, $p(G-v)\geq2>1$)。
\bu{推论3: 若二分图$G=(V_1,V_2,E)$有H回路, 则$|V_1|=|V_2|$。} (注意: 逆不成立, 这只是必要)

\rd{充分条件}: \bu{Ore道路定理(2.4.1): 若简单图$G$中任两个不相邻点$u,v$恒有$d(u)+d(v)\geq n-1$, 则$G$中存在H道路。}
\bu{Ore回路定理(推论2.4.1): 若简单图$G$中任两点$u,v$(不相邻)恒有$d(u)+d(v)\geq n$, 则$G$中存在H回路。}
\bu{Dirac定理(推论2.4.2): 若简单图$G$中每个节点$d(v)\geq n/2$, 则$G$中存在H回路。}
(注意: Ore/Dirac是充分非必要; 有些H图既不满足Ore也不满足Dirac)

\xt{闭合图与判定}: 引理2.4.1: 简单图$G$中$v_i,v_j$不相邻且$d(v_i)+d(v_j)\geq n$, 则$G$有H回路$\Leftrightarrow G+(v_i,v_j)$有H回路。\xt{闭合图$C(G)$}: 不断对满足$d(u)+d(v)\geq n$的不相邻点加边, 直到无法再加。\bu{引理2.4.2: $C(G)$唯一}。\bu{定理2.4.2: $G$存在H回路$\Leftrightarrow C(G)$存在H回路。}推论2.4.3: $C(G)=K_n\Rightarrow G$有H回路。

判定流程: 先试必要(割点/二分图不等$\Rightarrow$否); 再试充分(Ore/Dirac$\Rightarrow$是); 都不满足则算闭合图$C(G)$据Th2.4.2判定。

\subsection{最短路径}
\bu{定理: 正权图中两点间的最短路径一定是初级道路}。
\bu{引理2.6.1: 正权图中任一条最短路长度大于其局部路径长度。}
\bu{引理2.6.2: 正权图中若$P(i)$是$v_1$到$v_i$的最短路, 且$v_j\in P(i)$, 则$P(j)$也是$v_1$到$v_j$的最短路。}

\xt{Dijkstra算法(非负权单源$O(n^2)$)}: 按路径长度递增次序产生最短路径。核心: $T$集合中最小$d$值点必已确定最短, 移入$S$后松弛其邻点。
\alg{初始化: $d[s]=0$, 其余$d[v]=\infty$; $S=\emptyset$ \\ while $|S| < n$: \\\ST 选$u\notin S$且$d[u]$最小; $S = S\cup\{u\}$ \\\ST for 每条边$(u,v)$: \\\ST\ST if $d[u]+w(u,v) < d[v]$: $d[v]=d[u]+w(u,v)$}
路径存储用前驱数组$Q$: 更新$d[v]$时设$Q[v]=u$, 最终以$Q$反向追踪得路径。

\xt{Bellman-Ford($O(nm)$,可负权)}: 对所有边松弛$n-1$轮; 第$n$轮仍变则有负圈。边权全1用BFS$O(m)$。

\subsection{关键路径(AOE网)}
\bu{引理2.7.1: 有向图$G$无有向回路, 则一定存在入度为0的点及出度为0的点。}
\bu{定理2.7.1: 若有向图$G$无有向回路, 则可将$G$的节点重新编号为$v_1',v_2',\dots,v_n'$, 使得对任意边$v_i'v_j'\in E$, 都有$i<j$。}(构造法: 反复取入度0点号)

\xt{拓扑排序(Kahn算法,$O(n+m)$)}:
\alg{1. 计算每个节点的入度 \\ 2. 入度为0者入队$Q$ \\ 3. while $Q\neq\emptyset$: \\\ST $u = Q$.pop(); 输出$u$ \\\ST for $u$的每条出边$(u,v)$: \\\ST\ST $v$的入度$-1$; 若$v$入度变为0, $v$入$Q$}

\xt{关键路径}: 源到汇的最长路径决定项目最短完成时间。(1)按拓扑序前推: 事件最早发生时间$v_e(j)=\max_i\{v_e(i)+w_{ij}\}$, $v_e$(源)=0。(2)逆拓扑序后推: 事件最晚发生时间$v_l(i)=\min_j\{v_l(j)-w_{ij}\}$, $v_l$(汇)$=v_e$(汇)。(3)活动$(i,j)$的余量$=v_l(j)-v_e(i)-w_{ij}$。余量为0的活动是关键活动, 串联形成关键路径。

\subsection{中国邮路}
\rd{定理2.8.1: $L$是无向连通图$G$最佳邮路的充要条件是: (1)$G$的每条边最多重复一次; (2)在$G$的任意一个回路上, 重复边的长度之和不超过该回路长度的一半。} 充分性证明所有满足条件的重复边集总长相等。

\xt{奇偶点图上作业法}: (1)找出所有奇度点, 必为偶数个; (2)奇点配对, 每对沿最短路加重复边使全图偶度化; (3)检查每个回路: 若重复边总长$>$回路半长, 翻转(原重复变不重复, 原不重复变重复); (4)重复(3)直到全部满足, 走欧拉回路即得邮路。有向中国邮路要求强连通且每点入=出度, 否则添"入出度差"的最小费用流。Edmonds最小权匹配算法(1973)可求奇点最优配对。

TSP(旅行商): 找过所有点一次的最小权回路, NP-hard。分支定界法: 取下界$=$每行最小两元素和/2, 不断分支剪枝。便宜算法(近似$O(n^3)$): 任选起点, 每次选离当前子回路最近的点插入(贪心局部最优)。

% ====== 树 ======
\section{树 $\bigstar$}
\xt{定义3.1.1}: 不含任何回路的连通图称为\xt{树$T$}。树中度为1的节点为\xt{叶}, 度$>1$为\xt{分支点(内节点)}。不含回路的不连通图为\xt{林}。

\rd{定理3.1.2(树六等价定义, $n$阶$m$边无向图$G$)}:
(1)$G$连通且无回路(树的定义) $\Leftrightarrow$
(2)$G$中任意两个顶点之间存在唯一的路径 $\Leftrightarrow$
(3)$G$中无回路且$m=n-1$ $\Leftrightarrow$
(4)$G$连通且$m=n-1$ $\Leftrightarrow$
(5)$G$连通且$G$中任何边均为桥 $\Leftrightarrow$
(6)$G$中没有回路, 但在任何两个不同顶点间加一条新边, 所得图中得到唯一一个含新边的初级回路。
\xt{树的重要性质}: 任何树至少有两片叶; $\sum d(v)=2m=2n-2$; 树是边数最少的连通图, 也是边数最多的无回路图。

\xt{支撑树(Spanning Tree)$T$}: 无向图$G$的生成子图且为树。$T$中的边为\xt{树枝}, 不在$T$中的边为\xt{弦}。余树$G[E\setminus E(T)]$不一定连通但含回路。生成法: \xt{破圈法}(每次去回路一边, 共去$m-n+1$次); \xt{避圈法}(每次选不构成回路的边, 共选$n-1$条)。

全部支撑树枚举: \xt{删边/收缩递归}$\tau(G)=\tau(G-e)+\tau(G/e)$(不含$e$的树数$+$必含$e$的树数); 基本回路法(每条弦对应一基本回路, 弦替换回路中树枝); Gabow-Myers(利用关联矩阵行列式); Read-Tarjan最优$O(\tau+n+m)$。

\subsection{最小生成树(MST)}
连通加权图$G=\langle V,E,w\rangle$, $w:E\to\mathbb R^+$。\xt{MST}: 权之和最小的支撑树。
\bu{切割定理}: 设$S$是$V$的任意非空真子集, 若$e$是跨越$S$与$V-S$的最小权边, 则存在一棵包含$e$的最小生成树。
\bu{定理3.6.3}: 设$V'\subset V$是赋权连通图的节点真子集, $e$是二端点分跨$V'$和$V-V'$的最短边, 则$G$中一定存在包含$e$的最短树。
MST唯一性: 若所有边权值互不相同则MST唯一; 等权边存在时MST可能不唯一但总权相同。

\xt{Kruskal($O(m\log m)$)}: 边按权排序, 依次取不构成回路的边(并查集), 直到$n-1$条。

\xt{Prim算法(点扩展, 邻接阵$O(n^2)$)}:
\alg{1. $V'=\{v_0\}$, $T=\emptyset$ (任选起点$v_0$) \\ 2. while $V'\neq V$: \\\ST 找跨越$V'$与$V-V'$的最小权边$(u,v)$($u\in V',v\notin V'$) \\\ST $V' = V'\cup\{v\}$; $T = T\cup\{(u,v)\}$ \\ 返回$T$}

\subsection{关联矩阵与秩 $\bigstar$}
有向图$G$无环关联矩阵$B_{n\times m}$: $b_{ij}=1$若$v_i$是$e_j$始点; $-1$若$v_i$是$e_j$终点; $0$不关联。性质: 每列恰一个1和一个$-1$, 列和为0; 第$i$行1的个数$=d^+(v_i)$, $-1$个数$=d^-(v_i)$; 重边对应相同列。

\bu{定理3.2.1: 有向图$G$关联矩阵$B$的秩 $\operatorname{ran}B<n$。}

\xt{引理:} 设$B$为有向连通图$G$的关联矩阵, $C$是$G$中一个回路, 则$C$中各边对应$B$的各列线性相关。(回路子图的关联矩阵列向量线性相关, 添加其余点的零行后仍线性相关)

\rd{定理3.2.3: 有向连通图$G$关联矩阵$B$的秩 $\operatorname{ran}B=n-1$。}

\xt{基本关联矩阵$B_k$}: 在关联矩阵$B$中划去任意点$v_k$对应的行, 得$(n-1)\times m$矩阵。定理3.2.4: $B_k$的秩为$n-1$。推论3.2.1: $n$个点的树的基本关联矩阵秩为$n-1$。定理3.2.5: 设$B_k$为基本关联矩阵, $C$是$G$中一个回路, 则$C$中各边对应$B_k$的各列线性相关。推论3.2.2: 若$G$的子图$H$含有回路, 则$H$的诸边对应的$B_k$各列线性相关。
\rd{定理3.2.6: 令$B_k$是有向连通图$G$的基本关联矩阵, 则$B_k$的任意$(n-1)$阶子阵行列式$\neq0$当且仅当其各列所对应的边构成$G$的一棵支撑树。}(必要性: 行列式非零$\Rightarrow$无回路$\Rightarrow$树; 充分性: 树的基本关联矩阵满秩)
定理3.2.2: 设$B_0$为$B(G)$的任意一个$k$阶子方阵, 则$|B_0|=\pm1$或$0$(对$k$归纳, 每列最多2个非零元)。

\subsection{支撑树计数(矩阵树定理) $\bigstar$}
\rd{定理3.3.2(矩阵树定理): 设$B_k$是有向连通图$G$的某一基本关联矩阵, 则$G$的不同支撑树的数目 $\tau(G)=\det(B_k B_k^T)$。}

\xt{Binet-Cauchy定理(3.3.1)}: $A_{m\times n},B_{n\times m}(m\leq n)$, $\det(AB)=\sum_i A_iB_i$, $A_i$取$A$的$m$列子行列式, $B_i$对应$m$行。

\xt{特例:} 完全图$K_n$: $\tau(K_n)=n^{n-2}$(Cayley公式); 完全二分图$K_{p,q}$: $\tau(K_{p,q})=p^{q-1}q^{p-1}$。

\xt{计算步骤:} 写$B\to$划一行得$B_k\to$算$B_kB_k^T$乘积矩阵$\to$取行列式。不含边$e$: 算$\tau(G-e)$。必含边$e$: $\tau(G)-\tau(G-e)$; 或收缩$e$两端点算$\tau(G/e)$。无向图: 各边加方向后一一对应。


\subsection{根树计数 $\bigstar$}
\xt{根树(外向树)}: 有向树, 存在一根节点$v_0$入度0, 其余节点入度均为1。
\rd{定理3.3.3: 设$B_k$是有向连通图$G$的基本关联矩阵, $B_k^-$为将$B_k$中全部1改为0后的矩阵, 则$G$中以$v_k$为根的根树数目为 $\det(B_k^-(B_k^-)^T)$。}


\subsection{回路矩阵 $\bigstar$}
\xt{完全回路矩阵$C_e$}: 全部初级回路, $c_{ij}=1$若$e_j\in C_i$同向; $-1$反向; $0$不在。最多$2^{m-n+1}-1$行, 不一定独立。
\xt{基本回路矩阵$C_f$}: 取定生成树$T$, 每条余树边(弦)对应一个基本回路(由该弦+树上的唯一路径构成), 回路方向取该余树边的方向。重排边序(余树边在前,树枝在后,且次序与回路一致)$\Rightarrow C_f=[I_{(m-n+1)},\ C_{f12}]$。

\rd{定理3.4.1: $BC_e^T=0$}(关联矩阵$\times$回路矩阵转置$=0$, 边次序一致)。
\rd{定理3.4.2: $\operatorname{ran}C_e=m-n+1$。}
\xt{回路矩阵$C$}: 由$m-n+1$个互相独立的回路组成; 基本回路矩阵$C_f$即回路矩阵; $C=PC_f$($P$非奇异); $BC^T=0$。
定理3.4.3: 回路矩阵$C$的任一$(m-n+1)$阶子阵行列式非0当且仅当这些列对应于$G$的某一棵余树。

\subsection{割集矩阵 $\bigstar$}
\xt{边割集$S$}: 图$G=\langle V,E\rangle$的边子集, 满足: (1)$G'=\langle V,E-S\rangle$比$G$连通分支数多1; (2)对$\forall S'\subsetneq S$, $G''=\langle V,E-S'\rangle$与$G$连通分支数相同(极小性)。有向割集可给方向。
\xt{完全割集矩阵$S_e$}: 全部割集为行; $s_{ij}=1$同向, $-1$反向, $0$不属。
\xt{基本割集$S_f$}: 取树$T$, 每条树枝$e_i\in T$对应一个基本割集(含$e_i$及某些余树边), 方向与$e_i$一致。重排(余树边前,树枝后)$\Rightarrow S_f=[S_{f11},\ I_{(n-1)}]$。
\rd{定理3.4.5: $S_eC_e^T=0$。}(割集与回路若有交必交偶数条边, 相邻两条方向配合使乘积和为零)
\rd{定理3.4.6: $\operatorname{ran}S_e=n-1$。}(由$S_f$秩$n-1$及Sylvester定理)
定理3.4.7: 割集矩阵$S$的任一$(n-1)$阶子阵行列式非0当且仅当这些列对应于$G$的某棵支撑树。

\subsection{$\bigstar\bigstar\bigstar$ 三大矩阵关系(必考!)}
\xt{核心关系方程}: $BC^T=0$ (关联与回路正交), $SC^T=0$ (割集与回路正交)。

设$B_k=[B_{11},\ B_{12}]$(列重排: 余树边在前$|$树枝边在后):
\begin{itemize}
\item \rd{由$B_k$求$C_f$}: $C_f=[I,\ C_{f12}]$, 其中 \rd{$C_{f12}=-B_{11}^T B_{12}^{-T}$}(由$B_kC_f^T=0$推导)
\item \rd{由$C_f$求$S_f$}: $S_f=[S_{f11},\ I]$, 其中 \rd{$S_{f11}=-C_{f12}^T$}(Th3.4.8割集-回路关系)
\item \rd{由$B_k$求$S_f$}: \rd{$S_{f11}=B_{11}^{-1}B_{12}$}(推论3.4.1)
\end{itemize}
\xt{四步速算}: (1)确定支撑树$T$, 将边序重排为[余树边 $|$ 树枝边]; (2)写出$B_k$, 按列序分块$[B_{11},B_{12}]$; (3)求$B_{12}^{-1}$; (4)套公式得$C_{f12}=-B_{11}^TB_{12}^{-T}$, $C_f=[I,C_{f12}]$, $S_f=[-C_{f12}^T,I]$。验

\subsection{Huffman树与编码}
$m$叉树: 除叶外每点$\leq m$个子节点。完全二叉树: 每个内点恰2子。最优二叉树: 带权叶$w_i$, 路径长$\ell_i$, 使$WPL=\sum w_i\ell_i$最小。
\xt{Huffman算法}:
\alg{所有叶节点各自成树, 权如$w_i$, 全部入最小堆 \\ while 堆中元素数$>1$: \\\ST 取堆顶两棵最小权树(权$w_x,w_y$)\\ \ST 合并成一棵新树, 新树根权$=w_x+w_y$ \\\ST 新树入堆 \\ 堆中最后剩下的树即为最优二叉树}
\bu{定理3.5.1: Huffman算法得到的二叉树是最优二叉树。}前缀码: 任一字符编码不是另一字符编码的前缀; Huffman编码是最优前缀码。$m$叉Huffman: 需补足虚叶使$(n-1)\bmod(m-1)=0$, 虚叶权为0。

% ====== 平面图 ======
\section{平面图与对偶图}
平面图: 可画在平面上且任何两条边不在非顶点处相交。可嵌入平面但未画出$\to$\xt{可平面图}; 已画出$\to$\xt{平面嵌入}。

\xt{面(域)$R$}: 由图的若干边构成的区域, 内部不含任何节点及边。外部无限面$R_0$; 内部有限面。面$R$的\xt{边界}: 包围该域的诸边; \xt{次数$\deg(R)$}: 边界长度, 割边被计算2次。\bu{定理4.1.3: 所有面的次数之和等于边数的两倍 $\sum\deg(R_i)=2m$}。(每条边或在两个面边界上各算1次, 或是割边在一个面算2次)

\rd{欧拉公式(定理4.1.1): 对任意连通平面图 $n-m+r=2$。} 推论4.1.1: $k$个连通分支 $n-m+r=k+1$。推论4.1.2: 一般平面图 $n-m+r\geq2$。
\xt{定理4.1.2:} 设简单连通平面图$G$没有割边, 且每个域的边界数至少为$t$, 则 $m\leq\frac{t(n-2)}{t-2}$。用此证$K_5(t=3)$: $10>9$非平面; $K_{3,3}$(无三角$t=4$): $9>8$非平面。

\xt{极大平面图}($n\geq3$): 任意不相邻点加边后为非平面图的简单平面图。性质: 连通; 每个面次数均为3; $3r=2m$。\rd{定理4.2.1: 极大平面图 $m=3n-6,\ r=2n-4$。}推论4.2.1: 简单连通平面图 $m\leq3n-6$。无$K_3$(无三角)则$m\leq2n-4$。\bu{定理4.2.2: 简单连通平面图最小度 $\delta\leq5$。}

平面性检测: 定理4.3.1 平面图子图也平面; 定理4.3.2 非平面图母图也非平面; 推论4.3.1 $K_n(n\geq5)$和$K_{3,n}(n\geq3)$非平面; 定理4.3.3 重边自环不影响平面性。

同胚: 通过反复插入或消去2度顶点后同构。边收缩: 删$e=(u,v)$并$u,v$为新点$w$, $w$关联除$e$外$u,v$的一切边。
\rd{Kuratowski定理1: $G$是平面图$\Leftrightarrow G$中不含与$K_5$或$K_{3,3}$同胚的子图。}
\rd{Kuratowski定理2: $G$是平面图$\Leftrightarrow G$中没有可收缩到$K_5$或$K_{3,3}$的子图。}

对偶图$G^*$: $G$的每个面$f_i$中放顶点$v_i^*$; $G$的边$e$若在面$f_i$与$f_j$公共边界上, 做$G^*$边$e^*(v_i^*,v_j^*)$穿过$e$。性质4.4.1: $G^*$唯一且连通; 性质4.4.2 $G^*$连通; 性质4.4.3 $m^*=m,\ n^*=r,\ r^*=n$; 性质4.4.4 $d_{G^*}(v_i^*)=\deg_G(R_i)$; 性质4.4.5 $G$初级回路$\Leftrightarrow G^*$中割集。\bu{定理4.4.1: $G$有对偶图$\Leftrightarrow G$为平面图}。自对偶: $G^*\cong G$, 轮图皆为自对偶图。

% ====== 着色 ======
\section{图的着色}
\xt{点色数$\gamma(G)$}: 相邻点不同色的最少颜色数。$k$-可着色$\Leftrightarrow$存在$V$的$k$个独立集划分。特例: 零图1; $K_n:n$; $K_n-e:n-1$; 偶圈2奇圈3; 奇轮3偶轮4(轮$W_n$, $n$为点数); 含边二分图$\gamma=2$; 树$(n\geq2)$色数$=2$。

\bu{定理4.5.1: 非空图$\gamma(G)=2\Leftrightarrow G$没有奇回路$\Leftrightarrow G$是二分图。}(充分性: 取林$T'$, $\gamma(T)=2$, 偶回路加边不改变色; 必要性: 奇回路至少3色)
\bu{定理4.5.2: $\gamma(G)\leq\Delta(G)+1$}(不含自环)。
\bu{定理4.5.3(Brooks): 连通图$G$既不是完全图$K_n(n\geq3)$也不是奇圈, 则$\gamma(G)\leq\Delta(G)$。}

Welch-Powell近似着色: (1)点按度降序; (2)用色1涂第一个未涂点, 然后涂所有与其不相邻的未涂点; (3)换色重复直到全着色。不一定得最优。

定理4.5.4(收缩-加边递推): 设$i,j$为简单图$G$不相邻两节点, 则$\gamma(G)=\min\{\gamma(G+ij),\ \gamma(G/ij)\}$。($i,j$着同色$\to$收缩; 着异色$\to$加边, 二者必居其一)

\xt{色数多项式$f(G,t)$}: 用至多$t$种颜色给$G$着色的方案数。\bu{定理4.5.5: $f(K_n,t)=t(t-1)(t-2)\cdots(t-n+1)$}。\bu{定理4.5.6: $f(T_n,t)=t(t-1)^{n-1}$}。定理4.5.7: 不相邻$i,j$, $f(G,t)=f(G+ij,t)+f(G/ij,t)$。反复应用可将任意图化为完全图的和。

\xt{边色数$\gamma'(G)$}: 相邻边不同色的最少颜色数。可化为线图的点着色。\rd{定理4.6.4(Vizing): 简单图$\Delta(G)\leq\gamma'(G)\leq\Delta(G)+1$}。边色数只能取$\Delta$或$\Delta+1$, 但具体判定未解决。定理4.6.5: 轮图$\gamma'(W_n)=\Delta(W_n)=n-1$($n\geq4$)。定理4.6.6: $K_n$边色数 $n$为奇$\Rightarrow n$; $n$为偶$\Rightarrow n-1$。

\xt{面着色}: 平面图面着色$\Leftrightarrow$对偶图点着色。\bu{五色定理: 任何平面图都是5-面可着色的。} 四色定理: 平面图4-面可着色(1976年Appel\&Haken计算机辅助证明)。定理4.5.3: 有H回路$\Rightarrow$4色成立。定理4.5.4: 3-正则平面图4面可着$\Rightarrow$任意平面图4面可着。

% ====== 匹配 ======
\section{匹配}
\xt{匹配$M$}: 图$G$中互不共顶点的边子集。$M$的边关联的顶点为\xt{饱和点}, 否则为\xt{非饱和点}。\xt{极大匹配}: 加入任一其他边不再是匹配。\xt{最大匹配}: 边数最多的匹配。\xt{完备匹配}: 图所有顶点都饱和。

\xt{交互道路}: 属于$M$与不属于$M$的边交替出现的道路。\xt{可增广道路}: 两端点为非饱和点的交互道路(奇数条边, 不属于$M$的边比$M$中的多一条)。
\rd{定理5.1.1(Berge,1957): $M$为$G$的最大匹配当且仅当$G$不含关于$M$的可增广路径。}(必要性: 有可增广路则$M$非最大; 充分性: 无反例则$M$与最大匹配边数相等)

\xt{匈牙利算法(二分图最大匹配$O(nm)$)}:
\alg{初始$M=\emptyset$ \\ for 每个$X$部未饱和顶点$x_0$: \\\ST 从$x_0$出发做DFS/BFS找交互路 \\\ST 若找到$Y$部未饱和点$y$: // 找到可增广路 \\\ST\ST 翻转路径上所有边的匹配状态(匹配变非匹配, 非匹配变匹配) \\\ST\ST $|M|$增加1 \\\ST 若搜索完毕找不到可增广路: 标记$x_0$为"不可扩大"}

\xt{完全匹配}: 二分图$G=(V_1,V_2,E)$, $|V_1|\leq|V_2|$, 若$|M|=|V_1|$称$M$为$V_1$到$V_2$的完全匹配; 若$|V_1|=|V_2|$且$M$为完全匹配, 称\xt{完美匹配}。
\rd{定理5.2.1(Hall,1935): 二分图$G=(V_1,V_2,E)$($|V_1|\leq|V_2|$)中存在$V_1$到$V_2$的完全匹配, 当且仅当$V_1$中任意$k$个顶点至少与$V_2$中的$k$个顶点相邻($k=1,\dots,|V_1|$)。} (此条件称\xt{相异性条件}, 充分性
定理5.2.2: 二分图最大匹配边数$=|X|-\max_{A\subseteq X}\{|A|-|\Gamma(A)|\}$。
\rd{定理5.2.3(K\"onig, 1931): 二分图中, 最大匹配的边数$r$等于最小点覆盖的顶点数$s$。}(匹配-覆盖对偶)

\xt{KM算法(赋权二分图最佳匹配$O(n^3)$)}:
\alg{对每行取$l(x_i)=\max_j w_{ij}$为行顶标; $l(y_j)=0$为列顶标 \\ while 当前相等子图(边满足$l(x_i)+l(y_j)=w_{ij}$)无完美匹配: \\\ST 在相等子图上运行匈牙利找最大匹配 \\\ST 对未盖行/列求出$d=\min_{i\notin\text{行覆盖},j\in\text{列覆盖}}(l(x_i)+l(y_j)-w_{ij})$ \\\ST 所有未覆盖行顶标$-d$; 所有已覆盖列顶标$+d$ \\ 权重和 $=\sum l(x_i)+\sum l(y_j)$ 即为最大权和}

最小完美匹配: 取$C$为足够大常数, $w'_{ij}=C-w_{ij}$, 求$w'$的最大完美匹配即得$w$的最小完美匹配。

% ====== 网络流 ======
\section{网络流 $\bigstar$}
\xt{网络$N$}: 无自环有向连通图, 满足(1)唯一源$s$(入度0); (2)唯一汇$t$(出度0); (3)每条边$(i,j)$有非负容量$c_{ij}$。多源多汇转单源单汇: 加超源$s_0$连接所有源(容量$\infty$), 加超汇$t_0$被所有汇连接(容量$\infty$)。

\xt{流$f$}: 每条边赋非负实数, 满足(1)\xt{容量约束}: $0\leq f_{ij}\leq c_{ij}$; (2)\xt{流量守恒}: 对任意中间点$i$, $\sum_j f_{ij}=\sum_k f_{ki}$(流入=流出)。总流量$W=\sum_i f_{si}=\sum_j f_{jt}$。零流($f_{ij}\equiv0$)总是可行流。

\xt{割切$(S,\overline S)$}: $s\in S$, $t\in\overline S=V\setminus S$。\xt{割容量$C(S,\overline S)=\sum_{i\in S,\ j\in\overline S}c_{ij}$}(仅计从$S$到$\overline S$的正向边, 反向边不计)。\bu{定理6.1.1: 对任一可行流$f$和任一切割$(S,\overline S)$, $W\leq C(S,\overline S)$, 即 $\max W\leq\min C$。}

\xt{增流路径}: 从$s$到$t$的无向路径。(1)\xt{前向边}: 与$s\to t$方向一致, 若$f<c$可增$\delta=c-f$; (2)\xt{后向边}: 与$s\to t$方向相反, 若$f>0$可减$\delta=f$。路径增量$\delta=\min\{各边可增/减量\}$。增流后仍可行流。

\rd{定理6.1.2(最大流-最小割定理, Ford\&Fulkerson 1956): 网络流图中, 最大流量等于最小割切的容量: $\max W=\min C(S,\overline S)$。}
定义$S=\{v\mid$从$s$存在增流路径可达$v\}$。若$t\in S$, 则存在$s\to t$增流路径, 与$f$最大矛盾。故$t\notin S$。考虑割$(S,\overline S)$: 对所有$S\to\overline S$边, 必$f=c$(否则前向可增); 对$\overline S\to S$边, 必$f=0$(否则后向可减)。故$W=\sum_{s\to}v f_{si}=\sum_{i\in S,j\in\overline S}f_{ij}-\sum_{j\in\overline S,i\in S}f_{ji}=\sum_{i\in S,j\in\overline S}c_{ij}=C(S,\overline S)$。即最大流$=$该割容量, 而任何割容量$\geq$最大流, 故该割为最小割。

\rd{最小割物理意义}: 网络中的"瓶颈"或"咽喉"位置, 决定了网络的最大通过能力。

\xt{Ford-Fulkerson标号算法}:
\alg{(0) 取初始可行流$f$(如零流) \\ (1) 标号过程: \\\ST 给$s$标号$(-,\infty)$ \\\ST while $t$未被标号: \\\ST\ST 对每个已标号点$v_i$, 检查所有关联的未标号点$v_j$: \\\ST\ST\ST (a) 前向边$(v_i,v_j)$, 若$f_{ij}<c_{ij}$, 标$v_j$为$(v_i^+,\ \min\{\delta(v_i),\ c_{ij}-f_{ij}\})$ \\\ST\ST\ST (b) 后向边$(v_j,v_i)$, 若$f_{ji}>0$, 标$v_j$为$(v_i^-,\ \min\{\delta(v_i),\ f_{ji}\})$ \\ (2) 增流过程: \\\ST 若$t$被标号: 从$t$出发, 利用标号回溯找到增流路径$P$ \\\ST 对$P$上每条前向边$(u,v)$: $f_{uv}\gets f_{uv}+\delta(t)$ \\\ST 对$P$上每条后向边$(v,u)$: $f_{vu}\gets f_{vu}-\delta(t)$ \\\ST 清除除$s$外所有标号, 转(1) \\\ST 若$t$不能被标号: 当前流即为最大流, 已标号点集$S$与未标点集$\overline S$构成最小割}

\rd{Edmonds-Karp算法($O(nm^2)$)}: 在FF算法基础上, 每次沿\xt{最短增流路径}(边数最少)增流, 用BFS代替DFS进行标号。最多$O(nm)$次增广, 每次BFS$O(m)$。

\xt{Dinic算法($O(n^2m)$)}: 用BFS构建分层网络(各点到$s$的最短距离), 在分层网络上用DFS找阻塞流, 重复直到$t$不可达。点容量处理: 将点$v$拆为入点$v'$和出点$v''$, 中间连容量$c_{v'v''}=C_v$的边。应用: 任务分配(二分图最大流), 逃生路线(网格容量约束)。

\xt{最小费用流}: 边的容量$c_{ij}$之外还有单位流量费用$a_{ij}\geq0$。求流量为$W_0$的流中费用最小者。若$W_0$取最大流量, 则是\xt{最小费用最大流}。
\xt{最短增流路算法}: 以费用为边权, 每次在残量网络上找$s\to t$费用最短路(含负权反向边), 沿路增流。残量图边费: 若$0\leq f<c$, 正向边费$a_{ij}$; 若$f=c$, 删正向; 若$f>0$, 反向边费$-a_{ij}$; 若$f=0$, 删反向。

% ====== 代数结构 ======
\section{代数结构基础(七)}
\xt{映射$f:S\to T$}: 单射(不同元映到不同象); 满射($f(S)=T$); 双射(既是单射又是满射)。合成$g\circ f$满足结合律, 不满足交换律。定理7.1.1: $I_B\circ f=f,\ f\circ I_A=f$。
\bu{定理7.1.2: $f$左可逆$\Leftrightarrow f$单射; $f$右可逆$\Leftrightarrow f$满射; $f$可逆$\Leftrightarrow f$双射。}(入要单, 出要满)
定理7.1.3: 若$g\circ f=I_A,\ f\circ h=I_B$, 则$g=h$, 即可逆映射的逆唯一。推论: 双射的合成双射, $(g\circ f)^{-1}=f^{-1}\circ g^{-1}$。

\xt{等价关系$\sim$}: 自反性+对称性+传递性。等价类$\bar a=\{x|x\sim a\}$。\bu{定理7.2.1: 等价类要么相等, 要么交集为空。}\bu{定理7.2.2: 等价关系确定$A$的一个划分(等价类族构成对$A$的划分)}。\bu{定理7.2.3: 集合$A$的一个划分确定$A$上的一个等价关系。}商集$A/{\sim}$。

\xt{代数系统}: 非空集合$A$和其上封闭的若干运算$f_1,\dots,f_s$组成$(A,f_1,\dots,f_s)$。结合律$\Rightarrow$广义结合律(定理7.3.1): $x^m\cdot x^n=x^{m+n},\ (x^m)^n=x^{mn}$。\bu{定理7.3.2: 若既有左单位元$e_L$又有右单位元$e_R$, 则$e_L=e_R=e$且唯一。}\bu{定理7.3.3: 若代数系统有单位元$e$且结合律成立, 则若$x$有左逆元又有右逆元, 则左右逆元相等且唯一, $(x^{-1})^{-1}=x$。}

\xt{同态与同构}: 设$f:(X,\cdot)\to(Y,*)$。若$f(a\cdot b)=f(a)*f(b)$, $f$为\xt{同态}。若$f$双射则为\xt{同构}, 记$X\cong Y$。定理7.4.1: 同态象$(f(X),*)$是$(Y,*)$的子代数。定义7.4.5: 单同态(单)+满同态(满)$\to$同构。\bu{定理7.4.2(满同态保持代数性质): 满同态$f$保持运算的交换律、结合律; 若$X$有单位元$e$, 则$Y$有单位元$f(e)$; 若$x\in X$有逆元$x^{-1}$, 则$f(x)$有逆元$f(x^{-1})$。}自同态(同态到自身); 自同构(同构到自身)。

% ====== 群(1) ======
\section{群(1) 八-1\&八-2}
代数系统层级: 封闭$\to$封结(半群)$\to$封结幺(幺群/含幺半群)$\to$封结幺逆(群)。

\xt{半群(定义8.1.1)}: 非空集+结合律的二元运算。$(Z_m,+_m)$是半群。
\xt{幺群(定义8.1.2)}: 有单位元$e$的半群, $(Z_m,+_m,0)$。
\xt{交换幺群(定义8.1.3)}: 满足交换律的幺群。$(R,+,0)$是交换幺群非循环; $(Z_m,+_m,0)$是交换幺群且循环(生成元1)。
\xt{循环幺群(定义8.1.4)}: 存在生成元$g$, 任意元可表为$g^m$($m\geq0$)。$(N,+,0)$非幺群(无逆); $(R,+,0)$非循环(无单一生成元)。

\bu{定理8.1.1: 若二元运算适合结合律, 则适合广义结合律(指数律)。}$a^na^m=a^{n+m}, (a^n)^m=a^{nm}$。若$a$可逆, $a^{-n}=(a^n)^{-1}$。
\bu{定理8.1.2: 循环幺群是可交换幺群。}(证: $\forall a=g^m,b=g^n,\ ab=g^{m+n}=g^ng^m=ba$)
\bu{定理8.1.3: 满同态保持半群/幺群性。} 即若$S$是半群(幺群), 则其满同态象$f(S)$也是半群(幺群)。推论: 半群/幺群/子半群的同态象仍是半群/幺群/子半群。

\xt{群(定义8.2.1/8.2.2)}: 非空集$G$上二元运算满足: (1)结合律; (2)存在单位元$e$; (3)每元有逆元。可记为$(G,\cdot,e)$。等价定义: 含幺半群中所有元素可逆。
\xt{Abel群}: 满足交换律的群。平凡群$=\{e\}$; 有限群$|G|<\infty$; 无限群。

\xt{群的等价判定}:
\bu{定理8.2.2: 半群$(G,\cdot)$若有左单位元$e$且每元有左逆元($\forall a\exists a^{-1}:a^{-1}a=e$), 则$G$是群。}(证: 左单位元也是右单位元; 左逆元也是右逆元。需注意$e$必须是固定的左单位元)
\bu{定理8.2.3: 半群$(G,\cdot)$满足$\forall a,b$, 方程$ax=b$和$ya=b$在$G$中均有解, 则$G$是群。}(证: 构造左单位元再由Th8.2.2)

\rd{消去律与群}: 群必满足消去律。\rd{有限半群满足消去律必为群}(证: $x_1G=G=Gx_1\Rightarrow$有单位$\Rightarrow$有逆元)。无限未必(反例$(Z^+,\times)$)。

\bu{定理8.2.1(群基本性质)}: (1)单位元唯一; (2)每元逆元唯一; (3)指数律成立; (4)若$ab=ba$, 则$(ab)^n=a^nb^n$。
\bu{定理8.2.4(逆元性质)}: $(a^{-1})^{-1}=a$; \rd{$(ab)^{-1}=b^{-1}a^{-1}$}(注意顺序反转)。

\xt{元素的阶$O(a)$(定义8.2.4)}: 使$a^k=e$的最小正整数$k$; 若无此$k$则$O(a)=\infty$。
\bu{定理8.2.5(阶的性质)}: (1)$a^k=e\Leftrightarrow O(a)\mid k$; (2)$O(a^{-1})=O(a)$; (3)$r\leq|G|$, $e,a,\dots,a^{r-1}$两两不同。$O(b^{-1}ab)=O(a)$; $O(ab)=O(ba)$。有限群元素阶均有限; 无限群可全有限阶($P(\mathbb N),\oplus$)或仅$e$有限($Z,+$)。

\xt{子群$H\leq G$(定义8.2.5)}: $H\subseteq G$, 在$G$的运算下也构成群。$\{e\},G$为平凡子群; $H\neq G$为真子群。
\bu{定理8.2.6(子群判定三条件)}: $H$对运算封闭; $e\in H$; $\forall a\in H, a^{-1}\in H$。
\rd{定理8.2.7(子群判定一条件): $G$的非空子集$H$是$G$的子群$\Leftrightarrow\forall a,b\in H,\ ab^{-1}\in H$。}(证: $a=a\Rightarrow e\in H$; $e,b\Rightarrow b^{-1}\in H$; $a,b^{-1}\Rightarrow ab\in H$)

\xt{子群性质}: $H_1\cap H_2$子群; $\langle a\rangle=\{a^k\}$子群; $H\leq G\Rightarrow HH=H$。群不能是两真子群之并(Klein四元群是三真子群之并的特例)。

% ====== 群(2) ======
\section{群(2) 八-3\&八-4}
\subsection{循环群}
\xt{循环群(定义8.3.1)}: $G=\langle a\rangle=\{a^k\mid k\in Z\}$, $a$为生成元。$O(a)=n$$\Rightarrow n$阶循环群; $O(a)=\infty$$\Rightarrow$无限阶循环群。
\rd{定理8.3.1(生成元个数)}: $O(a)=\infty\Rightarrow$仅$a$和$a^{-1}$为生成元; $O(a)=n\Rightarrow\varphi(n)$个生成元($\varphi$为欧拉函数)。$a^k$为生成元$\Leftrightarrow\gcd(k,n)=1$。裴蜀定理: $\gcd(x,y)=1\Leftrightarrow\exists p,q:px+qy=1$。

\rd{定理8.3.2(循环群的子群)}: 循环群的子群仍是循环群。若$G=\langle a\rangle$无限, 非平凡子群也无限循环。若$G=\langle a\rangle$为$n$阶, 且$a^k$是子群$H$中$a$的最小正幂, 则$|H|=n/k$。\rd{对$n$的每个正因子$d$, 恰存在一个$d$阶子群$\langle a^{n/d}\rangle$; $n$阶循环群的子群个数$=n$的正因子个数。}


\xt{群同构(定义8.3.2)}: $f:G\to G'$双射且$f(ab)=f(a)f(b)$, 则$G\cong G'$。
\rd{定理8.3.4(循环群同构分类)}: $O(a)=\infty\Rightarrow G\cong(Z,+)$; $O(a)=n\Rightarrow G\cong(Z_n,+)$。推论: 任何两个同阶有限循环群同构。
\rd{定理8.3.5: 设$G$是群, $G'$是代数系统, 若存在双射$f:G\to G'$保持运算, 则$G'$也是群。}(用同构映射将群的性质"搬运"到新代数系统)

\subsection{变换群与置换群}
\xt{变换}: 集合$A$到自身的映射$f:A\to A$。一一变换$=$双射变换。$E(A)$是所有一一变换对乘法构成的群(一一变换群), 子群为\xt{变换群}。
有限集$|A|=n$上的一一变换为\xt{$n$元置换}。$S_n$: $n!$个置换, 对置换乘法构成\xt{$n$次对称群}。置换群$=S_n$的子群。

\xt{轮换$(i_1i_2\ldots i_l)$}: $\sigma(i_k)=i_{k+1},\sigma(i_l)=i_1$。对换$(i_1i_2)$为$l=2$。恒等置换$(i)$。
\bu{定理8.4.1: 不相交的轮换乘法可交换}(\text{证: 不相交意味着各自作用的元素互不干扰})。
\rd{定理8.4.2: $S_n$中任一置换可唯一分解为不相交轮换的乘积}(不计轮换次序)。

\bu{引理8.4.1(轮换的对换分解)}: $(i_1i_2\cdots i_k)=(i_1i_2)(i_2i_3)\cdots(i_{k-1}i_k)$ $=(i_1i_k)(i_1i_{k-1})\cdots(i_1i_2)$($k-1$个对换, 表示不唯一)。

\xt{逆序数$N(\sigma)$}: 排列中逆序总数。引理8.4.2: 对换个数的奇偶$=$逆序数的奇偶。奇置换$N$奇; 偶置换$N$偶。性质: 偶$\times$偶$=$偶, 奇$\times$奇$=$偶, 奇$\times$偶$=$奇。$\sigma$偶$\Leftrightarrow\sigma^{-1}$偶$\Leftrightarrow\tau^{-1}\sigma\tau$偶。轮换$(i_1\dots i_k)$阶$=k$, 是偶置换$\Leftrightarrow k$为奇($\because$可分解为$k-1$个对换)。
\rd{定理8.4.3(交错群$A_n$)}: $S_n$中所有偶置换构成子群$A_n$, $n\geq2$时$|A_n|=n!/2$。

\rd{定理8.4.4(Cayley定理)}: 任意群$G$与一个变换群同构, 特别地, $n$阶有限群与$S_n$的某个子群同构。证: 构造$\forall a\in G,\ f_a:x\mapsto ax$, 则$\mathbf G=\{f_a:a\in G\}$是变换群, 且$a\mapsto f_a$是$G\cong\mathbf G$的同构映射。

% ====== 群(3) ★★★ ======
\section{群(3) $\bigstar\bigstar\bigstar$ 八-5\&八-6}
\subsection{陪集与Lagrange定理}
$H\leq G$诱导的等价关系: $aRb\Leftrightarrow ab^{-1}\in H$(验证自反/对称/传递)。
\xt{左陪集$aH=\{ah\mid h\in H\}$}; \xt{右陪集$Ha=\{ha\mid h\in H\}$}。一般$aH\neq Ha$(如$S_3$中$H=\{e,(12)\}$时$(13)H\neq H(13)$)。

\rd{定理8.5.1(左陪集六条性质)}: 设$H\leq G$:
(1)$H=eH$, $a\in aH$;
(2)$|aH|=|H|$(由消去律保证);
(3)$a\in H\Leftrightarrow aH=H$;
(4)$\forall x\in aH$, $xH=aH$(陪集中任一元素可做代表元);
(5)$aH=bH\Leftrightarrow b^{-1}a\in H$;
(6)$aH\cap bH\neq\emptyset\Rightarrow aH=bH$, 即不同左陪集两两不相交。\rd{核心: $H$的所有左陪集构成$G$的一个划分!}

\rd{定理8.5.2(陪集分解)}: 有限群$G=a_1H\cup a_2H\cup\cdots\cup a_kH$(不相交划分), $k=[G:H]$为$H$在$G$中的\xt{指数}(陪集个数)。取$H=\{e\}$则$k=|G|$, $[G:\{e\}]=|G|$。

\rd{Lagrange定理: 对有限群$G$及其子群$H$, $|G|=[G:H]\cdot|H|$, 即$|H|$整除$|G|$。}(证: $G$被划分为$[G:H]$个大小等于$|H|$的陪集)

\rd{推论1(元素阶整除群阶)}: 设$|G|=n$, $\forall a\in G$, $O(a)\mid n$, 且$a^n=e$。(证: $\langle a\rangle$是子群, $|\langle a\rangle|=O(a)$, 由Lagrange $O(a)\mid n$)
\rd{推论2(素数阶群)}: 阶为素数$p$的群必为循环群。(证: 取$a\neq e$, $O(a)$整除$p$, 只能$O(a)=p$, 故$G=\langle a\rangle$)
\rd{推论3(子群乘积的阶)}: $|AB|=\dfrac{|A|\cdot|B|}{|A\cap B|}$, 其中$AB=\{ab\mid a\in A,b\in B\}$。(证: $AB=\bigcup_{a\in A}aB$, 不同陪集数$=[A:A\cap B]$)

\xt{典型例题(6阶群必有3阶子群)}: 设$|G|=6$。非$e$元阶2,3,6。6阶元$\Rightarrow\langle a^2\rangle$3阶; 3阶元直接; 全2阶$\Rightarrow$交换$\Rightarrow$两2阶元生4阶子群($4\nmid6$矛盾)。故必存3阶子群。

\subsection{正规子群与商群 $\bigstar$}
\rd{定义8.6.1(正规子群)}: $H\trianglelefteq G$当且仅当$\forall a\in G,\ aH=Ha$(左右陪集相等)。$\{e\},G$恒正规; 交换群的任何子群皆正规。

\rd{定理8.6.1(正规等价条件, $H\leq G$)}: 以下等价:
(1)$H\trianglelefteq G$ $\Leftrightarrow$ (2)$\forall g\in G,\ gHg^{-1}=H$ $\Leftrightarrow$ (3)$\forall g\in G,\ gHg^{-1}\subseteq H$ $\Leftrightarrow$ \bu{(4)$\forall g\in G,\forall h\in H,\ ghg^{-1}\in H$(最实用判定!)}。
(证: (1)$\Rightarrow$(2): $gHg^{-1}=(gH)g^{-1}=(Hg)g^{-1}=H$; (2)$\Rightarrow$(3)显然; (3)$\Rightarrow$(4)显然; (4)$\Rightarrow$(1): $gH\subseteq Hg$且$Hg\subseteq gH$故相等)

\rd{定理8.6.2(正规子群运算)}: (1)若$A,B\trianglelefteq G$, 则$A\cap B\trianglelefteq G$, $AB\trianglelefteq G$。(2)若$A\trianglelefteq G$且$B\leq G$, 则$A\cap B\trianglelefteq B$, $AB\leq G$(注意仅$AB$是子群, 不一定是正规子群)。
口诀: 正规$\cap$正规$=$正规(对$G$); 正规$\times$正规$=$正规(对$G$); 正规$\cap$普通$=$正规(对普通); 正规$\times$普通$=$普通(仅子群, 不保证正规)。

\rd{定理8.6.3(商群)}: 设$H\trianglelefteq G$。$G/H=\{gH\mid g\in G\}$在\xt{陪集乘法}$(aH)(bH)=(ab)H$下构成群, 称为$G$关于$H$的商群。证:(1)封闭性利用$H\trianglelefteq G$保证; (2)结合律继承自主群; (3)单位元$eH=H$; (4)逆元$(aH)^{-1}=a^{-1}H$。$|G/H|=[G:H]$。

\xt{例}: $Z/mZ\cong Z_m$; 换位子群$G'$($xyx^{-1}y^{-1}$生成)$\trianglelefteq G$, $G/G'$交换。

\section*{矩阵计算四步}
求$C_f,S_f$: (1)定树$T$, 列重排[余树$|$树枝]; (2)$B_k=[B_{11},B_{12}]$; (3)求$B_{12}^{-1}$; (4)$C_{f12}=-B_{11}^TB_{12}^{-T}\Rightarrow C_f=[I,C_{f12}],\ S_f=[-C_{f12}^T,I]$。\\
$\tau(G)=\det(B_kB_k^T)$; 根树$=\det(B_k^-(B_k^-)^T)$(1改0)。验证 ran$B_k=n-1$, ran$C_f=m-n+1$, ran$S_f=n-1$。

\end{multicols*}
\end{document}
